% ECONOMICS_PROBLEM: {"id": "C31-442", "title": "Interference and equilibrium", "jel_code": "C31", "source_id": "OP-0421"}
% !TeX program = xelatex
\documentclass[11pt,a4paper]{article}
\usepackage[margin=23mm]{geometry}
\usepackage{fontspec}
\setmainfont{Latin Modern Roman}
\usepackage{amsmath,amssymb,mathtools}
\usepackage{xcolor,enumitem,booktabs,longtable,array,xurl,hyperref}
\definecolor{muted}{HTML}{53636C}
\hypersetup{colorlinks=true,urlcolor=blue,linkcolor=blue}
\setlength{\parindent}{0pt}
\setlength{\parskip}{5pt}
\newcommand{\R}{\mathbb R}
\newcommand{\E}{\mathbb E}
\newcommand{\N}{\mathbb N}
\newcommand{\ind}{\mathbf 1}
\newcommand{\Var}{\operatorname{Var}}
\newcommand{\Cov}{\operatorname{Cov}}
\newcommand{\supp}{\operatorname{supp}}
\DeclareMathOperator*{\argmin}{arg\,min}
\DeclareMathOperator*{\argmax}{arg\,max}
\newcommand{\entrymeta}[3]{{\sffamily\small\color{muted}\textbf{\texttt{#1}}\quad|\quad #2\par #3\par}}
\newcommand{\Problemref}[1]{\texttt{#1}}
\newcommand{\JELref}[2]{\texttt{#2} (JEL \texttt{#1})}
\begin{document}
\section*{Interference and equilibrium}
\textbf{Problem ID:} \texttt{C31-442}\quad \textbf{JEL:} \texttt{C31}\par
% BEGIN ECONOMICS STATEMENT
\label{problem:OP-0421}
{\sffamily\small\color{muted}Original subject group: Econometrics, causal inference and measurement\par}
\label{definitions:problem:30007703}
\textbf{Audit classification:} Published research agenda. The publication explicitly motivates interference and equilibrium-adjustment research. The revised finite-population estimand distinguishes assignment distributions from full-policy endpoints.
\subsection*{Definitions and precise research target}
Consider a finite population of product markets, each containing sellers and buyers whose prices, entry and purchase choices adjust after a platform policy. For market \(m\), let \(Z_m\) be the vector of seller treatment assignments, and let \(Y_m(Z_m)\) be total monetary consumer and producer surplus after a prespecified adjustment horizon, with producer surplus measured net of real costs. Arbitrary within-market interference is allowed; markets are assumed independent conditional on recorded common shocks. The target is
\[
\tau=E[Y_m(\mathbf1)-Y_m(\mathbf0)],
\]
where \(\mathbf1\) and \(\mathbf0\) assign the policy to every or no eligible seller. An experiment may randomize only intermediate treatment saturations, so it directly identifies outcomes for the tested assignment distributions rather than the two endpoints. Determine which additional full-saturation market experiments, exposure restrictions or equilibrium models permit informative identification or bounds for \(\tau\). Specify any exposure mapping and restrictions on price adjustment before estimation, and distinguish immediate assignment effects from outcomes after entry and prices settle. Construct uncertainty statements that reflect market-level randomization and untested saturation extrapolation. The agenda explicitly recognizes interference and equilibrium adjustment; the particular surplus measure, horizon and proposed endpoint estimand here are an adaptation rather than a theorem already stated by the source.

\textbf{Definitions, assumptions and mathematical qualifications.} Let the fixed market cohort have weights \(\nu_m\geq0\), \(\sum_m\nu_m=1\), so \(\tau=\sum_m\nu_m\{Y_m(\mathbf1)-Y_m(\mathbf0)\}\); superpopulation expectation is a separate target. Each potential outcome includes the complete market assignment and a fixed response horizon, with common monetary welfare units and equilibrium selection. An exposure mapping is a given function \(g_m(Z_m)\) and imposes \(Y_m(Z_m)=\widetilde Y_m(g_m(Z_m))\); it is a substantive restriction. Intermediate-saturation randomization identifies averages over its tested assignment law. Endpoint identification requires positive endpoint assignment probability or specified extrapolation assumptions. Conditional independence of market shocks does not remove within-market spillovers.
\subsection*{Variants and assumption changes}
\begin{enumerate}
\item \textbf{Editorial, endpoint design.} Randomize entire markets to all-treated or all-control, estimate the finite-cohort full-policy contrast with market-level randomization uncertainty.
\item \textbf{Editorial, untested endpoints.} Test only partial saturations, impose a stated Lipschitz or monotonicity restriction on selected equilibrium welfare as a function of saturation, and bound the two endpoints; test the restriction where possible.
\end{enumerate}
\subsection*{References and source audit}
\textbf{Checked source.} 2025 conference draft, Section 1.1, third research direction, PDF p. 5. Full primary text retrieved. \url{https://conference.nber.org/confer/2025/DTs25/farronato.pdf}.
\begin{enumerate}\raggedright
\item\hypertarget{definitions:source:30007703:1}{} Chiara Farronato. 2025. \emph{Digital Platforms as Information Aggregators: Implications for their Design and Regulation}. 2025 conference draft, Section 1.1, third research direction, PDF p. 5.
\url{https://conference.nber.org/confer/2025/DTs25/farronato.pdf}.
\end{enumerate}

\subsection*{Common conventions: Source mathematical conventions}

These conventions apply to each entry unless explicitly replaced.

\textbf{Models and identification.} A primitive model \(m\) specifies all probability laws, feasible choices, information, equilibrium and measurement rules used by its target. The admissible class \(\mathcal M\) is fixed independently of outcomes. If \(P_O(m)\) is its observable probability law and \(\tau(m)\) the target, its identified set at observable law \(P\) is
\[
 \mathcal I_\tau(P)=\{\tau(m):m\in\mathcal M,\ P_O(m)=P\}.
\]
Point identification means that this set is a singleton. An empty set signals incompatibility of data and assumptions. Coordinate bounds need not describe the joint set. Bounds are sharp only if every retained target is attainable or arbitrarily approachable by compatible models. Finite bounds require explicit restrictions on unobserved quantities; unrestricted confounding, hidden wealth, extrapolation or arbitrary equilibrium selection can yield uninformative sets.

\textbf{Causal interventions.} A potential outcome \(Y_i(p)\) is the outcome under a complete policy regime \(p\), including timing, financing, information and market spillovers. The target population and units are fixed before treatment. With interference, use the complete assignment vector or a justified exposure mapping; individual no-interference notation alone cannot identify an economy-wide intervention. A difference of mean potential outcomes does not require the cross-world joint distribution. Conditioning on post-treatment migration, survival, enrollment or employment changes the estimand and needs separate justification.

\textbf{Statistical statements.} An estimator, confidence set, forecast or test specifies a sampling law, model class and loss/coverage criterion. Uniform \((1-\alpha)\)-coverage means \(\inf_{m\in\mathcal M}\Pr_m\{\tau(m)\in C_n\}\geq1-\alpha\), or an explicitly asymptotic version. The whole parameter space trivially has coverage, so informativeness requires a stated width, risk or power objective. A forecasting improvement is relative to a named benchmark, information set and evaluation population.

\textbf{Equilibrium and optimization.} An equilibrium comprises feasible plans, optimality, beliefs where needed, budget constraints and all market/resource clearing. The primitives or a stated benchmark define each correspondence \(\mathcal E(p;m)\). Nonemptiness is assumed or proved, never inferred just from compact policy sets. A selected-equilibrium target uses a declared selection rule; a robust target quantifies over all permitted equilibria. Use suprema or infima unless attainment is established. An \(\varepsilon\)-optimal policy is feasible and within \(\varepsilon\) of the finite optimum value. Report an empty feasible set separately.

\textbf{Welfare and accounting.} Normative utility, interpersonal normalization, social weights, numeraire, discounting and the use of public receipts are primitives. Behavioral utility need not coincide with the welfare criterion. Household welfare includes ownership income and rents once; adding profits again to compensating variation generally double-counts them. A monetary equivalent is defined by an explicit transfer and price convention, with monotonicity and range conditions. Average effects, mechanism contributions, transfers and welfare losses are different targets. Infinite sums require convergence and dynamic feasible plans require terminal or no-Ponzi conditions.

\textbf{Source versions.} A working-paper date, revision date and journal publication year are distinguished. A DOI for a journal version is not automatically the DOI of an earlier manuscript. Locators refer to the stated version and printed pages where specified. A metadata-only check does not verify a claim about a full-text conclusion. Every entry's source subsection narrows the provenance claim accordingly.
% END ECONOMICS STATEMENT
\end{document}
